%% paper2_revised.tex - revised manuscript after the ICACS-2026 rejection.
%% Venue-neutral article format with author-year citations (IJPHM uses APA
%% style; swap in the IJPHM template once the target is confirmed).
%% Every number in this file comes from results/revised/*.json; see
%% revision/08_change_log.md for what changed and why.
\documentclass[11pt]{article}
\usepackage[T1]{fontenc}
\usepackage[margin=1in]{geometry}
\usepackage{amsmath,amssymb}
\usepackage{graphicx}
\usepackage{booktabs}
\usepackage{xcolor}
\usepackage{url}
\usepackage[round,authoryear]{natbib}
\usepackage[hidelinks]{hyperref}
\graphicspath{{figures/}{./}}

\title{Protocol Sensitivity in Unsupervised Industrial Anomaly Detection:
Threshold Transfer, Retraining Risk, and Inference Cost}
\author{Sarjan Shah \and Abhi Shah\\[2pt]
\small Department of Information Technology, Lok Jagruti University, Vadodara, Gujarat, India\\
\small \texttt{shahsarjan968@gmail.com}, \texttt{abhi.shah@ljku.edu.in}}
\date{}

\begin{document}
\maketitle
\let\thefootnote\relax\footnotetext{Sarjan Shah et al. This is an open-access article distributed under the terms of the Creative Commons Attribution 3.0 United States License, which permits unrestricted use, distribution, and reproduction in any medium, provided the original author and source are credited.}

\begin{abstract}
Comparisons of unsupervised anomaly detectors for industrial sensor data often
report a winning method, but the winner can depend on evaluation choices that are
rarely varied. This paper shows how three such choices change the conclusions:
which data define the normal training period (a widely shared re-upload of SWaT
places the attack period's normal rows in the file named for normal operation),
how detection thresholds are calibrated, and whether retraining data are selected
with labels. Nine lightweight detectors and a small LSTM autoencoder are
evaluated under standard protocols with threshold-free metrics (AUROC, AUPRC),
threshold-dependent metrics under two calibration rules, event recall and
block-bootstrap confidence intervals on four benchmarks: SKAB (pump and valve
testbed), SWaT and WADI (water treatment and distribution, official iTrust releases)
and HAI (turbine, boiler and water-treatment hardware-in-the-loop testbed). On SWaT,
all ten detectors rank attacks equally well (AUROC $0.81$--$0.86$, overlapping
intervals); on SKAB with pooled training none is reliably better than chance, while
SKAB's per-run protocol yields run-level AUROC up to $0.79$ and an inverted ranking;
on HAI, PCA reconstruction and LOF lead ($0.93$, $0.92$) and Elliptic Envelope trails
($0.66$). Protocol choices moved MCC by up to $0.8$ for the same detector. Thresholds
set only from normal data failed wherever chemical analysers drifted between the
normal and the attack recordings: on SWaT and WADI, distance- and reconstruction-based
detectors raised alarms on 42--75\% of normal readings, while detectors that score each
sensor's rarity separately stayed below 12\%. Adding one drifting analyser (AIT201,
missing from a public re-upload of SWaT) raised the $z$-score's false-alarm rate from
1.4\% to 42\%; removing three drifting analysers from WADI cut false alarms from 56\%
to about 1\%. In a streaming
simulation, SWaT's normal operation drifted enough to raise a static Isolation
Forest's false-alarm rate from 12\% to 76\%; retraining restored detection only when
retraining data were known to be normal; a drift trigger saved some refits relative to
a schedule but its checks cost more than refitting every hour; and label-free
retraining failed. Single-reading
cost in a one-CPU, 1\,GB container ranged from microseconds to about 700\,ms. The
results argue for choosing detectors by threshold stability and retraining needs and
for reporting protocols in enough detail to tell which of these a benchmark number
measures.
\end{abstract}

\noindent\textbf{Keywords:} unsupervised anomaly detection; multivariate time
series; evaluation protocol; concept drift; retraining policy; resource-constrained inference

\section{Introduction}

Pumps, valves and treatment plants are instrumented with sensors whose joint
behaviour changes before or during a fault. Because faults and attacks are rare
and costly to label, condition monitoring often relies on \emph{unsupervised}
(novelty) detectors that are fitted to normal operation and flag departures from
it \citep{chandola2009,mukherjee2021}. In practice, false alarms decide whether such a
monitor is trusted, which is why deployed systems add contextual filtering and
adaptive thresholds on top of the detector \citep{lee2017}. For an engineer choosing such a detector for a
resource-limited monitoring node, three practical questions matter: which
detector separates anomalies from normal operation best, whether that answer
survives the passage of time once the detector is deployed, and what the
detector costs to run.

Published comparisons give less guidance on these questions than their
headline numbers suggest. Large benchmark studies report that no family of
methods is consistently best and that deep models often do not beat
well-chosen classical ones \citep{audibert2022,schmidl2022,sarfraz2024}. A
parallel line of work shows that evaluation choices, such as point adjustment,
metrics that reward trivial detectors, and flawed benchmark labels, can
manufacture apparent progress \citep{kim2022,garg2021,wu2023,liu2024}. These
critiques concern mainly metrics and labels. Less attention has been paid to
the other protocol choices an applied study must make: how the training and
test data are defined, how a detection threshold is calibrated, and, once
deployed, how and on which data a detector is retrained.

This study began as a conventional detector comparison on the Skoltech Anomaly
Benchmark (SKAB) \citep{skab} and the Secure Water Treatment testbed (SWaT)
\citep{goh2016,mathur2016}. An initial analysis, built with choices that are
common in applied work, named clear winners on each dataset and a large benefit
from drift-triggered retraining. A subsequent audit showed that those conclusions
rested on three protocol choices: the SWaT copy used (a public re-upload) stores
the normal-labelled rows of the attack period in the file named for normal
operation, so the model had been trained on the normal intervals \emph{between}
attacks rather than on the separate normal run; thresholds had been calibrated
with different rules for different detectors; and retraining data had been
selected with ground-truth labels. Instead of only correcting these choices, this
paper measures how much each one changes the conclusions, re-runs SWaT on the official
release, and adds two further benchmarks (HAI and WADI) to test whether the corrected
findings hold.

The contributions are empirical and methodological, not algorithmic:
\begin{enumerate}
\item A controlled comparison of nine lightweight unsupervised detectors and a
small LSTM autoencoder on SKAB, SWaT, WADI and HAI under standard, documented protocols,
reported with threshold-free metrics (AUROC, AUPRC), threshold-dependent
metrics (MCC, event recall, false-alarm rate) and block-bootstrap confidence
intervals (Section~\ref{sec:results}).
\item A protocol-sensitivity analysis showing how the split definition, the
training protocol, the threshold rule and the sensor set change the measured
performance and the ranking of detectors, including a check of a public SWaT re-upload
against the official release (Section~\ref{sec:protocol}).
\item A retraining-policy study on SWaT's genuine normal run and attack period
that compares never, periodic, drift-triggered and continuous retraining, with
retraining data selected with or without labels (Section~\ref{sec:drift}).
\item Single-reading latency, throughput, model size and peak memory for every
detector inside a one-CPU, 1\,GB ARM64 container (Section~\ref{sec:edge}).
\end{enumerate}

The study answers four research questions:
\begin{description}
\item[RQ1] Under standard, documented protocols with uncertainty estimates, do
lightweight detectors differ in how well they rank anomalies on SKAB, SWaT, WADI and HAI?
\item[RQ2] How much do the split definition, the training protocol and the
threshold rule change measured performance and the ranking of detectors?
\item[RQ3] Does drift-triggered retraining outperform periodic or continuous
retraining, and what happens when retraining data are selected without labels?
\item[RQ4] What does each detector cost per reading under a one-CPU, 1\,GB
memory cap?
\end{description}

\section{Related Work}
\label{sec:related}

\paragraph{Detectors.} Classical unsupervised detectors differ in what they
treat as normal. Distance- and density-based methods such as the Local Outlier
Factor (LOF) score a point by the relative density of its neighbourhood
\citep{breunig2000}; boundary methods such as the One-Class SVM enclose the
bulk of the training data \citep{scholkopf2001}; covariance methods flag large
Mahalanobis distances, with the Minimum Covariance Determinant (MCD) giving a
robust estimate \citep{rousseeuw1999}; PCA reconstruction scores the residual
outside a low-rank normal subspace \citep{shyu2003}; Isolation Forest isolates
points by random partitioning \citep{liu2008}; and the histogram- and
copula-based HBOS, COPOD and ECOD combine per-feature tail probabilities and are
fast and nearly parameter-free \citep{goldstein2016,li2022ecod,zhao2019}. Deep
sequence models, including LSTM encoder-decoders \citep{malhotra2016},
prediction-error detectors \citep{hundman2018}, adversarial autoencoders
\citep{audibert2020}, graph networks \citep{deng2021} and transformers
\citep{tuli2022}, report strong results on SWaT, usually with protocols and
metrics that differ from paper to paper.

\paragraph{What benchmark studies agree on.} Comparisons that put conventional,
machine-learning and deep detectors side by side find that no family wins
consistently \citep{audibert2022}, that multivariate problems are harder on
average \citep{schmidl2022}, and that simple or linear baselines match many
recent deep models once inflated metrics are removed \citep{sarfraz2024}. Tabular
benchmarks reach similar conclusions \citep{goldstein2016,han2022}. A result of
the form ``detector X wins on dataset A but not on B'' is therefore expected
and is not, on its own, a contribution.

\paragraph{Where evaluations go wrong.} Point adjustment credits a whole event
when any of its points is flagged and can make random scores look competitive
\citep{kim2022}. Point-wise F1 and accuracy can reward detectors that flag
almost everything, which the Matthews correlation coefficient (MCC) does not
\citep{chicco2020,garg2021}. Threshold-dependent metrics make comparisons hinge
on the threshold rule; threshold-free AUC measures and their range-aware
extensions avoid this \citep{paparrizos2022,liu2024}, and range-based
precision and recall account for events rather than points \citep{tatbul2018}.
Benchmark data can itself be flawed \citep{wu2023}, and the PHM community has
called for shared measurement and evaluation practice for health-monitoring methods
\citep{weiss2021}. These studies examine metrics
and labels; the protocol choices examined here (split semantics, threshold
calibration rule, and the label assumption in retraining) are complementary.

\paragraph{Drift and retraining.} Concept-drift research distinguishes
detecting a change, understanding it and adapting to it \citep{gama2014,lu2018}.
Two-sample tests on feature distributions are a simple and competitive way to
detect shift \citep{rabanser2019}; adaptive windows such as ADWIN resize the
reference automatically \citep{bifet2007}; streaming anomaly benchmarks such as
NAB score detectors online \citep{ahmad2017}. For an anomaly detector, adapting
to drift carries a specific risk: data used for retraining may contain the very
anomalies the detector should flag. Retraining policies have been compared
directly for anomaly detectors in IT operations, where periodic (blind)
retraining performed slightly better than retraining triggered by a drift
detector and the best retraining window depended on the detector
\citep{poenaru2024}; a recent preprint on streaming classifiers likewise reports
periodic retraining outperforming error- and ADWIN-triggered policies except under
recurring drift \citep{dasari2026}. Those studies retrain on data whose labels or
normality are known. The present study asks the same question for lightweight
detectors on industrial-control data and adds the factor a deployed detector cannot
avoid: whether the retraining data are selected with labels, by the detector's own
predictions, or not at all.

\paragraph{Resource cost.} MLPerf Tiny measures accuracy, latency and energy of
single inferences on constrained hardware, including an anomaly-detection task
\citep{banbury2021}. Comparable cost reporting for classical multivariate
detectors on industrial data is rare. Here, cost is measured as single-reading
latency, batch throughput, model size, load time and peak memory under explicit
CPU and memory caps. Energy is not measured.

\section{Datasets}
\label{sec:data}

\paragraph{SKAB.} SKAB was recorded on a water-circulation testbed (pump,
pipes, valves) with eight sensors sampled at 1\,Hz: two accelerometers (RMS),
motor current, pressure, water temperature, pump-housing temperature, voltage
and volumetric flow rate \citep{skab}. It contains one 9{,}401-row anomaly-free
run and 34 fault runs (46{,}860 rows in total) in three families: valve~1
(16 runs), valve~2 (4) and other faults (14). Faults are physically induced and
labelled per second. Each file is an independent experiment.

\paragraph{SWaT.} SWaT is a six-stage water-treatment plant with 51 sensors and
actuators sampled at 1\,Hz \citep{goh2016,mathur2016}. The official release used here
(iTrust, SUTD: SWaT.A1 \& A2, December 2015) has a seven-day normal run
(\texttt{Normal\_v1}, 22--28 December, 495{,}000 rows, the first 30 minutes of start-up
already removed) and an attack period (\texttt{Attack\_v0}, 28 December--2 January,
449{,}919 rows) in which 54{,}621 rows (12.1\%) belong to 36 attacks (37 rows spelled
"A ttack" in the label column are counted as attacks). Following common practice
\citep{li2019madgan}, the normal run is used from 22 December 22:00, six hours after the
plant started. The 25 continuous sensors (flows, levels, pressures and chemical
analysers) are used; the 26 discrete actuator states are not.

A widely downloaded public re-upload of the same data (Kaggle,
\texttt{vishala28/swat-dataset-secure-water-treatment-system}, v3) was used in the
initial analysis. Its \texttt{normal.csv} holds the normal run twice, with seven
channels blank (including the analyser AIT201), plus the normal-labelled rows of the
attack period, and its \texttt{attack.csv} holds only the attack rows. Rebuilding the
standard split from it by de-duplicating timestamps reproduces the official row counts;
Section~\ref{sec:protocol} compares results from the two copies.

\paragraph{WADI.} WADI is a water-distribution testbed built as an extension of SWaT,
with 127 points sampled at 1\,Hz \citep{ahmed2017}. The official A2 release (iTrust,
November 2019) provides a 14-day normal run (784{,}571 rows, 25 September--7 October
2017) and a two-day attack file (172{,}803 rows, 9--11 October) in which 9{,}977 rows
(5.8\%) belong to 14 attacks. Four columns are empty in both files and are dropped; the
first 21{,}600 normal rows are removed as for SWaT; the 67 points with more than ten
distinct values in the training rows are used. The file's time column is malformed, so
rows are indexed by their 1\,Hz order.

\paragraph{HAI.} HAI (version 21.03) comes from a hardware-in-the-loop testbed in
which a steam-turbine process, a boiler process and a water-treatment process are
coupled through a simulator, sampled at 1\,Hz \citep{shin2020,shin2021}. It provides
three normal training files (921{,}603 rows) and five test files (402{,}005 rows) in
which 2.2\% of rows belong to 50 scripted attacks. Of its 78 points, the 55 with more
than ten distinct values in the training files are treated as continuous and used.
HAI adds a second process type beyond water and a test set with a much lower
anomaly rate than SWaT.

\section{Method}
\label{sec:method}

\subsection{Features}
For each continuous sensor, the raw value, 5-second and 30-second causal
rolling means, a 5-second rolling standard deviation and the first difference
are computed (40 features on SKAB, 125 on SWaT, 335 on WADI, 275 on HAI). Features are computed on each
complete recording \emph{before} any rows are removed, so windows never join
readings across a removed fault and, at test time, use only past readings.
The initial analysis' elapsed-time feature was removed: its range differs between
training (up to 9{,}960\,s on SKAB) and test segments (at most 291\,s), and it
has no meaning for a monitor that runs continuously. Its effect is reported
in Section~\ref{sec:protocol}. Wall-clock features remain excluded because they
leak the identity of the recording session.

\subsection{Detectors}
Nine classical detectors are evaluated, all producing a continuous anomaly score:
a maximum absolute $z$-score; PCA reconstruction error (95\% of variance);
Elliptic Envelope (robust MCD on ten whitened principal components, needed
because the features are collinear); Isolation Forest (100 trees); One-Class SVM
(RBF kernel, $\nu{=}0.1$); LOF ($k{=}20$, novelty mode); and ECOD, COPOD and
HBOS from PyOD \citep{zhao2019}, with library defaults. On SWaT, WADI and HAI,
One-Class SVM and LOF are fitted on a 20{,}000-row evenly spaced subsample of the training data
because their cost grows quadratically; on WADI and HAI, ECOD and COPOD are fitted on a
100{,}000-row subsample because PyOD rescores test data jointly with the stored
training data and the full set exceeded 16\,GB of memory (test data are scored in
50{,}000-row chunks, so their WADI and HAI scores depend slightly on chunk size). All other
detectors use every training row. The deep baseline is an LSTM autoencoder (one-layer encoder and decoder,
64 hidden units, $\approx$\,52k--58k parameters) on 30-second windows of raw
scaled sensors, trained for 25 epochs with Adam; a row's score is the
reconstruction error of the window ending at that row \citep{malhotra2016}.
All models are fitted on rows labelled normal, with one exception: under the
SKAB reference protocol the first 400 rows of each run are used as given, and
2.1\% of them are labelled anomalous. Table~\ref{tab:impl} lists the implementation
details.

\begin{table}[t]
\centering\small
\caption{Implementation details. All classical detectors receive the same standardised
feature matrix (StandardScaler fitted on training rows only); scores are oriented so
that higher means more anomalous.}
\label{tab:impl}
\begin{tabular}{lll}
\toprule
Detector & Configuration & Score (higher = more anomalous) \\
\midrule
$z$-score & training mean and s.d.\ per feature & $\max_j |x_j-\mu_j|/\sigma_j$ \\
PCA-recon & components for 95\% of variance & squared reconstruction error \\
Elliptic Envelope & 10 whitened PCs; robust MCD (default support) & Mahalanobis distance \\
Isolation Forest & 100 trees, default subsample, seeds 0--9 & $-$score\_samples \\
One-Class SVM & RBF, $\nu{=}0.1$, gamma = scale & $-$decision\_function \\
LOF & $k{=}20$, novelty mode & $-$score\_samples \\
ECOD, COPOD, HBOS & PyOD defaults & decision\_function \\
LSTM-AE & 1-layer LSTM enc./dec., 64 units, window 30, Adam, lr $10^{-3}$, & window MSE \\
 & batch 256, 25 epochs, stride 2 (4 on HAI); raw sensors, own scaler & \\
\midrule
\multicolumn{3}{l}{Software: Python 3.14 (host) / 3.12 (container), NumPy 2.4.6, SciPy 1.17.1, pandas 3.0.3,}\\
\multicolumn{3}{l}{scikit-learn 1.9.0, PyOD 3.6.7, PyTorch 2.12.1. Seeds: 0 unless stated; bootstrap: NumPy default\_rng(0).}\\
\bottomrule
\end{tabular}
\end{table}

\subsection{Threshold rules}
\label{sec:thr}
A score becomes an alarm when it exceeds a threshold $\tau$. Because the
comparison should not depend on one rule, two rules are applied identically to
every detector:
\begin{itemize}
\item \textbf{Label-free:} $\tau$ is the 99th percentile of scores on held-out
\emph{normal} data, so about 1\% of normal readings are expected to raise an
alarm. No anomaly labels are used.
\item \textbf{Validation-MCC:} $\tau$ maximises MCC on a labelled validation
set that contains anomalies, excluding thresholds that flag more than 99\% of
rows. This is semi-supervised calibration and is reported as such.
\end{itemize}
A third rule, maximising validation F1, is used only to show how much the
initial analysis' choice of F1 mattered.

\subsection{Splits}
\paragraph{SKAB, pooled protocol.} As in the initial analysis, each fault run is
split chronologically into its first 80\% and last 20\%. The training portions of
20\% of the runs (whole runs, fixed seed) form a labelled validation set; the
model is trained on the anomaly-free run plus the normal rows of the remaining
training portions; the last 20\% of every run is the test set (7{,}503 rows,
30.0\% anomalous).
\paragraph{SKAB, benchmark protocol.} The SKAB reference implementation trains
a model per run on its first 400 rows and tests on the rest. This protocol is
also reported, with the label-free threshold set from each run's own training
scores, so results can be read against the SKAB leaderboard.
\paragraph{SWaT, standard protocol.} Detectors are trained on the normal run
from 22 December 22:00 to 27 December 13:18 (400{,}680 rows); the rest of the run
(74{,}520 rows) is held-out normal data
for the label-free threshold; the first third of the attack period is the
labelled validation set for the validation-MCC threshold. Threshold-free metrics
and label-free results are reported on the full attack period; both threshold
rules are also compared on its last two thirds, which neither rule has seen.
\paragraph{WADI.} The same design is used on WADI: training on the first 85\% of the
normal run, the last 15\% as held-out normal data, the first third of the attack file
for the validation-MCC threshold, and metrics on the whole file and on its last two
thirds.
\paragraph{HAI.} The same design is used: training on the three normal files except
the last 15\% of the third, which is held-out normal data for the label-free
threshold; the first test file is the labelled validation set for the
validation-MCC threshold; metrics are reported on all five test files and, for the
rule comparison, on test files 2--5. Bootstrap blocks are one-hour blocks within
each test file.

\subsection{Metrics and uncertainty}
AUROC and average precision (AUPRC) summarise ranking quality independent of
any threshold. At a threshold, precision, recall, F1, MCC, false-alarm rate on
normal rows, event recall (the fraction of contiguous anomalous segments with at
least one alarm) and median detection delay are reported. MCC is
\begin{equation}
\mathrm{MCC}=\frac{TP\cdot TN-FP\cdot FN}{\sqrt{(TP+FP)(TP+FN)(TN+FP)(TN+FN)}}
\end{equation}
which is zero for any predictor independent of the labels, including one that
flags every row; F1 for that predictor equals $2\pi/(1+\pi)$ for anomaly
prevalence $\pi$ (0.46 on the SKAB test set), which is why F1 alone can credit
a detector that separates nothing \citep{chicco2020}. Precision is bounded by
prevalence in the same way: with false-alarm rate $f$ and recall $r$,
$\mathrm{precision}=\pi r/(\pi r+(1-\pi)f)$, so the same detector looks much
worse in precision on SWaT ($\pi{=}0.121$) than on a test set built with
$\pi{=}0.315$, as in the initial analysis.
Test rows are strongly autocorrelated, so 95\% confidence intervals are
obtained by resampling whole blocks with replacement: the 34 test files on SKAB
and non-overlapping one-hour blocks on SWaT. The same resamples are used for
every detector, so paired differences are valid. Isolation Forest is repeated
over ten seeds and the LSTM autoencoder over three seeds on SKAB.

\subsection{Drift and retraining}
The retraining study simulates deployment on SWaT. An initial model is fitted
on the first 50{,}000 normal rows after start-up (about 14 hours). The
remainder of the normal week and then the attack period are streamed in
one-hour batches; each batch is scored \emph{before} it can be used for
retraining. Four policies are compared: \emph{static} (never retrain);
\emph{drift-triggered} (retrain when the mean two-sample Kolmogorov--Smirnov
statistic between the batch and the current reference,
$\bar D=\frac{1}{d}\sum_{j=1}^{d}\sup_x\left|F^{\mathrm{ref}}_j(x)-F^{\mathrm{win}}_j(x)\right|$
\citep{massey1951}, exceeds the 95th percentile of its within-initial-window
values); \emph{periodic} (the same number of retrains as the triggered policy,
evenly spaced); and \emph{always} (after every batch). Each retrain uses the
last five batches, with rows selected in one of three ways: \emph{oracle}
(ground-truth normal rows, as in the initial analysis), \emph{predicted-normal}
(rows the current model did not flag) or \emph{all} rows. After every fit the
threshold is the 99th percentile of the fit-window scores. Isolation Forest and
Elliptic Envelope are tested. The scaler is fitted once on the initial window.
The KS statistic is used instead of its $p$-value because, with thousands of
rows per batch, $p$-values vanish for negligible differences, whereas $\bar D$
is a bounded effect size.

\subsection{Resource measurement}
Each detector is fitted and timed in its own freshly started (spawned) process
inside a Docker container run with \texttt{--cpus=1 --memory=1g} on an ARM64
(Apple M4) host, with BLAS threads limited to one; the container's control group
confirmed a 1.0-CPU quota and a 1{,}024\,MB memory limit. Reported are fit time,
serialised model size, load time from the serialised form, single-reading latency
over 1{,}000 calls after one warm-up call (median, 95th and 99th percentiles), batch
throughput on 10{,}000 rows, and memory in two parts: the process's resident memory
after the interpreter, libraries and data are loaded (\emph{baseline}), and the peak
increase above that baseline while fitting and scoring (\emph{model memory}, from
the kernel's high-water mark reset after loading). The container-wide peak and
memory events are read from the control group.
The cost of recomputing the rolling features for one new reading from a
30-row buffer is reported separately. This is a capped container, not an edge
device: a physical device's slower cores, thermal limits and I/O are not
represented, so latencies are lower bounds.

\section{Results}
\label{sec:results}

\subsection{SKAB under the pooled protocol}
Table~\ref{tab:skab} and Fig.~\ref{fig:auroc} (left) report SKAB under the pooled
protocol. Every detector is close to chance. The highest AUROC is Isolation
Forest's $0.575$ ($95\%$ CI $0.504$--$0.645$; mean $0.575\pm0.018$ over ten seeds),
and its CI barely excludes $0.5$; the CIs of the top five detectors overlap, and
the paired difference between the two best (Isolation Forest and COPOD) is
$0.012$ (CI $-0.066$ to $0.092$). The LSTM autoencoder is below chance
(AUROC $0.442$, three seeds within $\pm0.001$). With a calibrated threshold the
best MCC is $0.113$ (COPOD; CI $0.045$--$0.190$). The label-free rule gives MCC
within $\pm0.05$ of zero for every detector: thresholds learned from normal
data in this pooled setting do not separate faults.

\begin{table}[t]
\centering\small
\caption{SKAB, pooled protocol (test: last 20\% of each of 34 runs, 7{,}503 rows,
30.0\% anomalous). AUROC and MCC with 95\% block-bootstrap CIs over files. MCC
columns use the label-free (LF) or validation-MCC (VM) threshold; event recall
uses VM; 33 fault events.}
\label{tab:skab}
\begin{tabular}{lcccccc}
\toprule
Detector & AUROC [95\% CI] & AUPRC & MCC LF & MCC VM [95\% CI] & Event rec. VM \\
\midrule
Isolation Forest & 0.575 [0.504, 0.645] & 0.338 & 0.005 & 0.097 [$-$0.026, 0.233] & 0.94 \\
COPOD & 0.558 [0.507, 0.621] & 0.320 & $-$0.042 & 0.113 [0.045, 0.190] & 1.00 \\
ECOD & 0.535 [0.476, 0.601] & 0.315 & $-$0.016 & 0.066 [$-$0.010, 0.159] & 1.00 \\
Elliptic Envelope & 0.520 [0.459, 0.581] & 0.305 & $-$0.050 & 0.051 [$-$0.079, 0.182] & 0.45 \\
$z$-score & 0.513 [0.477, 0.554] & 0.307 & $-$0.051 & $-$0.032 [$-$0.102, 0.039] & 0.48 \\
HBOS & 0.510 [0.420, 0.602] & 0.315 & 0.041 & 0.012 [$-$0.116, 0.152] & 0.85 \\
One-Class SVM & 0.498 [0.463, 0.533] & 0.300 & $-$0.035 & $-$0.042 [$-$0.106, 0.009] & 0.73 \\
PCA-recon & 0.489 [0.455, 0.517] & 0.294 & $-$0.046 & $-$0.047 [$-$0.104, $-$0.005] & 0.76 \\
LOF & 0.489 [0.450, 0.533] & 0.299 & 0.007 & $-$0.021 [$-$0.073, 0.032] & 0.82 \\
LSTM-AE & 0.442 [0.392, 0.500] & 0.267 & $-$0.075 & $-$0.078 [$-$0.165, $-$0.007] & 0.27 \\
\bottomrule
\end{tabular}
\end{table}

\begin{figure}[t]
\centering
\includegraphics[width=\textwidth]{auroc_ci.png}
\caption{AUROC with 95\% block-bootstrap confidence intervals (blocks: test files on
SKAB, one-hour blocks on SWaT, WADI and HAI). Dashed line: chance. On SKAB (pooled) no
detector is clearly above chance; on SWaT (official release) all ten are well above
chance and statistically indistinguishable; on WADI only the per-feature detectors are
above chance while three analysers drift; on HAI a leading group separates from HBOS and
Elliptic Envelope.}
\label{fig:auroc}
\end{figure}

\subsection{SWaT under the standard protocol}
Table~\ref{tab:swat} and Fig.~\ref{fig:auroc} report SWaT (official release, 25
sensors). All ten detectors rank attack rows well (AUROC $0.811$--$0.856$) and their
CIs overlap almost completely; the difference between the top two (ECOD and COPOD) is
$0.007$ (CI $-0.003$ to $0.017$). Isolation Forest's AUROC over ten seeds is
$0.819\pm0.005$. The LSTM autoencoder ranks attacks as well as the classical detectors
(AUROC $0.822$, AUPRC $0.731$).

The detectors separate once a threshold is applied. With the label-free rule, which
targets 1\% false alarms on held-out normal-week data, COPOD reaches MCC $0.628$ (CI
$0.488$--$0.730$), Isolation Forest $0.539$ and HBOS $0.463$, with false-alarm rates of
0.1--11.8\% on the attack period's normal rows. The $z$-score, One-Class SVM, LOF, PCA
reconstruction and the LSTM autoencoder instead raise alarms on 42--75\% of the attack
period's \emph{normal} rows, although their thresholds gave exactly 1\% on held-out
normal-week data. Their scores shift between the normal week and the attack week
independently of the attacks. The shift comes from SWaT's chemical analysers: relative
to the training range, AIT402 is out of range on 98\% of the attack period's normal
rows, AIT201 on 73\% and AIT501--504 on 53--58\%. Detectors that combine all features
through distances or reconstructions inherit any single feature's shift; detectors that
score each feature's rarity separately (COPOD, ECOD, HBOS) or split on one feature at a
time (Isolation Forest) bound its influence. Elliptic Envelope shows the opposite failure,
with very few alarms (MCC $0.005$). When a labelled third of the attack period is used to
calibrate thresholds, MCC on the remaining two thirds rises sharply for the detectors
whose label-free thresholds failed (One-Class SVM $0.822$, $z$-score $0.820$, LOF
$0.742$), but not for PCA reconstruction ($0.078$) or the LSTM autoencoder ($0.101$).

\begin{table}[t]
\centering\small
\caption{SWaT, official release, standard protocol, 25 sensors (test: full attack
period, 449{,}919 rows, 12.1\% attacks, 35 contiguous attack segments). AUROC and MCC
with 95\% block-bootstrap CIs (one-hour blocks). LF: label-free threshold (99th
percentile of held-out normal-week scores); FAR: false-alarm rate on the attack
period's normal rows; VM: validation-MCC threshold calibrated on the first third,
scored on the last two thirds.}
\label{tab:swat}
\begin{tabular}{lcccccc}
\toprule
Detector & AUROC [95\% CI] & AUPRC & MCC LF [95\% CI] & FAR LF & Event rec. LF & MCC VM \\
\midrule
ECOD & 0.856 [0.777, 0.925] & 0.694 & 0.336 [0.216, 0.425] & 0.002 & 0.26 & 0.617 \\
COPOD & 0.849 [0.770, 0.918] & 0.726 & 0.628 [0.488, 0.730] & 0.001 & 0.17 & 0.788 \\
HBOS & 0.826 [0.735, 0.903] & 0.553 & 0.463 [0.306, 0.584] & 0.118 & 0.63 & 0.440 \\
Isolation Forest & 0.824 [0.726, 0.895] & 0.563 & 0.539 [0.369, 0.675] & 0.070 & 0.69 & 0.552 \\
LSTM-AE & 0.822 [0.722, 0.899] & 0.731 & 0.130 [0.056, 0.200] & 0.742 & 0.89 & 0.101 \\
Elliptic Envelope & 0.821 [0.727, 0.890] & 0.509 & 0.005 [$-$0.031, 0.052] & 0.015 & 0.74 & 0.401 \\
LOF & 0.820 [0.721, 0.897] & 0.718 & 0.149 [0.090, 0.212] & 0.748 & 1.00 & 0.742 \\
One-Class SVM & 0.816 [0.711, 0.896] & 0.718 & 0.184 [0.106, 0.274] & 0.581 & 0.91 & 0.822 \\
PCA-recon & 0.814 [0.715, 0.893] & 0.722 & 0.131 [0.055, 0.208] & 0.726 & 0.91 & 0.078 \\
$z$-score & 0.811 [0.708, 0.890] & 0.719 & 0.215 [0.109, 0.331] & 0.421 & 0.91 & 0.820 \\
\bottomrule
\end{tabular}
\end{table}

\subsection{HAI}
\label{sec:hai}
Table~\ref{tab:hai} reports HAI. Unlike SWaT, the detectors do separate. PCA
reconstruction (AUROC $0.929$, CI $0.899$--$0.954$), LOF ($0.920$) and the LSTM
autoencoder ($0.916$) lead, followed by One-Class SVM and the $z$-score ($0.895$);
the difference between the top two is $0.008$ (CI $-0.009$ to $0.025$). HBOS
($0.728$) and Elliptic Envelope ($0.663$, CI $0.608$--$0.716$) are clearly worse,
and Isolation Forest sits in between ($0.837$; $0.815\pm0.009$ over ten seeds).
Label-free thresholds transfer well on HAI: held-out normal data from the training
period predicted the test-period false-alarm rate closely (0.4--1.6\% for every
classical detector against a 1\% target; 9.1\% for the LSTM autoencoder), and the
leading five detectors raised alarms in all 50 attacks. With these thresholds MCC is
$0.597$--$0.627$ for PCA reconstruction, LOF, One-Class SVM and the $z$-score.
Calibrating on a labelled test file raises MCC on the remaining files only slightly
(by $0.05$--$0.06$) for the leaders, but by $0.33$ for the LSTM autoencoder, whose
label-free threshold was too low.

Across the two plants, the detectors that lead differ: PCA reconstruction and LOF,
which failed under label-free thresholds on SWaT, lead on HAI, where no analyser drift of
the SWaT kind occurs and thresholds transfer; Elliptic Envelope trails on both.

\begin{table}[t]
\centering\small
\caption{HAI 21.03 (test: five files, 402{,}005 rows, 2.2\% attacks, 50 attacks).
AUROC and label-free MCC with 95\% block-bootstrap CIs (one-hour blocks). FAR:
false-alarm rate on normal test rows; VM: validation-MCC threshold calibrated on test
file 1, scored on files 2--5.}
\label{tab:hai}
\begin{tabular}{lcccccc}
\toprule
Detector & AUROC [95\% CI] & AUPRC & MCC LF [95\% CI] & FAR LF & Event rec. LF & MCC VM \\
\midrule
PCA-recon & 0.929 [0.899, 0.954] & 0.641 & 0.624 [0.572, 0.670] & 0.016 & 1.00 & 0.678 \\
LOF & 0.920 [0.895, 0.948] & 0.617 & 0.626 [0.570, 0.675] & 0.016 & 1.00 & 0.691 \\
LSTM-AE & 0.916 [0.887, 0.948] & 0.585 & 0.336 [0.284, 0.391] & 0.091 & 1.00 & 0.664 \\
One-Class SVM & 0.895 [0.858, 0.930] & 0.590 & 0.627 [0.573, 0.677] & 0.012 & 1.00 & 0.662 \\
$z$-score & 0.895 [0.860, 0.928] & 0.539 & 0.597 [0.544, 0.648] & 0.013 & 1.00 & 0.653 \\
ECOD & 0.865 [0.817, 0.906] & 0.470 & 0.519 [0.446, 0.580] & 0.007 & 0.78 & 0.515 \\
COPOD & 0.853 [0.795, 0.900] & 0.479 & 0.531 [0.467, 0.591] & 0.004 & 0.80 & 0.513 \\
Isolation Forest & 0.837 [0.783, 0.884] & 0.359 & 0.414 [0.343, 0.484] & 0.004 & 0.66 & 0.425 \\
HBOS & 0.728 [0.652, 0.792] & 0.115 & 0.145 [0.078, 0.214] & 0.011 & 0.42 & 0.121 \\
Elliptic Envelope & 0.663 [0.608, 0.716] & 0.114 & 0.179 [0.119, 0.248] & 0.007 & 0.48 & 0.177 \\
\bottomrule
\end{tabular}
\end{table}

\subsection{WADI}
\label{sec:wadi}
WADI is the hardest benchmark (Table~\ref{tab:wadi}). With all 67 points, the
per-feature detectors lead (HBOS AUROC $0.741$, Isolation Forest $0.720$, COPOD and ECOD
$0.71$; Isolation Forest $0.725\pm0.016$ over ten seeds), while the $z$-score, One-Class
SVM, LOF, PCA reconstruction, Elliptic Envelope and the LSTM autoencoder are at chance
(AUROC $0.46$--$0.52$). These six detectors raise alarms on 55.8--56.1\% of the attack
file's normal rows under the label-free rule, and the near-identical rate points to a
single cause: the analyser 2B\_AIT\_002\_PV is outside its training range on 55.5\% of
those rows (2A\_AIT\_004\_PV on 26\%, 1\_AIT\_003\_PV on 7\%). Repeating the experiment
without these three analysers removes the effect: every detector's label-free false-alarm
rate falls to 0.6--1.4\%, AUROC moves into a common band of $0.650$--$0.748$ with
overlapping CIs, and label-free MCC reaches $0.351$ for the LSTM autoencoder and
$0.18$--$0.23$ for LOF, the $z$-score, PCA reconstruction and One-Class SVM. Absolute
performance on WADI remains low, as other studies also find.

\begin{table}[t]
\centering\small
\caption{WADI A2 (test: 172{,}803 rows, 5.8\% attacks, 14 attacks), with all 67 points
and without the three drifting analysers (64 points). AUROC with 95\% block-bootstrap CI
(one-hour blocks); label-free MCC and false-alarm rate on normal test rows.}
\label{tab:wadi}
\begin{tabular}{lcccccc}
\toprule
 & \multicolumn{3}{c}{All 67 points} & \multicolumn{3}{c}{Without 3 drifting analysers} \\
Detector & AUROC [95\% CI] & MCC LF & FAR LF & AUROC [95\% CI] & MCC LF & FAR LF \\
\midrule
HBOS & 0.741 [0.588, 0.863] & 0.078 & 0.014 & 0.748 [0.599, 0.865] & 0.017 & 0.013 \\
Isolation Forest & 0.720 [0.555, 0.843] & 0.031 & 0.013 & 0.720 [0.570, 0.835] & 0.029 & 0.013 \\
COPOD & 0.710 [0.582, 0.818] & 0.055 & 0.009 & 0.709 [0.576, 0.817] & 0.054 & 0.009 \\
ECOD & 0.709 [0.585, 0.810] & 0.061 & 0.009 & 0.703 [0.567, 0.808] & 0.060 & 0.009 \\
One-Class SVM & 0.517 [0.380, 0.655] & $-$0.030 & 0.559 & 0.733 [0.590, 0.859] & 0.183 & 0.011 \\
LOF & 0.484 [0.337, 0.626] & $-$0.026 & 0.559 & 0.650 [0.502, 0.763] & 0.230 & 0.006 \\
$z$-score & 0.479 [0.341, 0.620] & $-$0.027 & 0.559 & 0.697 [0.575, 0.807] & 0.214 & 0.009 \\
PCA-recon & 0.478 [0.330, 0.627] & $-$0.027 & 0.558 & 0.686 [0.480, 0.833] & 0.192 & 0.010 \\
Elliptic Envelope & 0.475 [0.327, 0.635] & $-$0.074 & 0.559 & 0.709 [0.551, 0.844] & $-$0.006 & 0.011 \\
LSTM-AE & 0.464 [0.295, 0.623] & 0.014 & 0.561 & 0.720 [0.494, 0.887] & 0.351 & 0.011 \\
\bottomrule
\end{tabular}
\end{table}

\subsection{How much protocol choices change the conclusions}
\label{sec:protocol}

\paragraph{Split definition (SWaT).} Scoring the same detectors with the same
code on the re-upload's layout (the split of the initial analysis: train on the
normal rows between attacks, test on 31.5\% attacks) gives AUROC $0.840$--$0.934$
and reorders the detectors: HBOS ranks first there and last under the standard
split with the same 24 sensors (Fig.~\ref{fig:protocol}, left). Under that layout the training data come
from the same five days as the test attacks, so the normal-week to attack-week
shift described above never appears in evaluation.

\paragraph{Data copy and sensor set (SWaT).} With the same 24 sensors, the standard
split rebuilt from the public re-upload and the official release give practically
identical results (every detector's AUROC within $0.001$, label-free MCC within $0.008$),
so the reconstruction is faithful and the re-upload's problem lies in how its files are
organised. The one real difference is the 25th sensor, the analyser AIT201, which the
re-upload leaves blank. Adding it leaves the ranking tied (AUROC $0.81$--$0.86$) but
changes which thresholds survive: the $z$-score's label-free false-alarm rate rises from
1.4\% to 42\% (MCC $0.716$ to $0.215$) and One-Class SVM's from 3.8\% to 58\% (MCC
$0.660$ to $0.184$), while COPOD (MCC $0.592$ to $0.628$) and Isolation Forest ($0.630$ to
$0.539$) are barely affected. A single sensor therefore decides whether the simplest
detector looks dependable.

\paragraph{Threshold rule (SWaT).} The initial analysis tuned one contamination
value on Isolation Forest and applied it to every model. On the same layout,
per-model validation-MCC thresholds (with the revised detector settings of
Section~\ref{sec:method}) raise the $z$-score's MCC from $0.180$ to
$0.844$, PCA reconstruction's from $0.002$ to $0.834$ and LOF's from $0.105$ to
$0.830$, while Elliptic Envelope, the initial winner at $0.766$, moves only to
$0.810$. The initial ranking therefore reflected mainly the shared threshold,
not the detectors.

\paragraph{Training protocol (SKAB).} Under the SKAB reference protocol (one
model per run, trained on its first 400 rows), the same detectors are far
better than under pooled training: mean per-run AUROC reaches $0.794$
(PCA reconstruction), $0.785$ (LOF) and $0.771$ (One-Class SVM), against
$0.489$--$0.575$ pooled (Fig.~\ref{fig:protocol}, right; Table~\ref{tab:skabbench}).
The ranking inverts: Isolation Forest, the best pooled detector, is among the
weakest per run, and ECOD and COPOD fall to chance because, with only 400 training
rows and a 99th-percentile threshold, they almost never alarm. Of the first 400
rows, 2.1\% are labelled anomalous, so this protocol trains on slightly
contaminated data, as in the reference implementation. Absolute numbers are not
directly comparable with the SKAB leaderboard (best F1 $0.78$ at a 13.6\% false-alarm
rate for a convolutional autoencoder), because features, thresholds and post-processing
differ; detectors here reach similar F1 ($0.72$--$0.73$) only at much higher
false-alarm rates (47--71\%), while Elliptic Envelope offers a more balanced
operating point (F1 $0.644$, false alarms 20.2\%, MCC $0.361$).

\paragraph{Threshold rule (SKAB).} On the pooled SKAB test set, Isolation
Forest's MCC is $0.005$ with the label-free rule, $0.097$ with validation-MCC and
$0.118$ with validation-F1; the F1 rule buys its MCC with recall $0.71$ at
higher false-alarm rates. Differences of this size are within the CIs.

\paragraph{Feature choices.} Removing the elapsed-time feature and computing
features before dropping fault rows changed SKAB AUROC by at most $0.02$
(Isolation Forest $0.591$ initial pipeline, $0.555$ with elapsed time, $0.575$
revised), well within the CIs; these corrections matter for validity, not for
the numbers.

\begin{figure}[t]
\centering
\includegraphics[width=\textwidth]{protocol_sensitivity.png}
\caption{Same detectors, same code, different protocol. Left: SWaT AUROC under the
re-upload layout used in the initial analysis versus the standard split. Right: SKAB
AUROC with one model trained on all runs (pooled) versus one model per run trained
on its first 400 rows (SKAB reference protocol; mean over runs).}
\label{fig:protocol}
\end{figure}

\begin{table}[t]
\centering\small
\caption{SKAB reference protocol (per run: first 400 rows train, remainder test;
label-free threshold from each run's training scores). FAR and MAR as defined by the
SKAB benchmark (false-alarm and missed-alarm rates).}
\label{tab:skabbench}
\begin{tabular}{lccccc}
\toprule
Detector & F1 & FAR & MAR & MCC & Mean run AUROC \\
\midrule
Elliptic Envelope & 0.644 & 0.202 & 0.445 & 0.361 & 0.749 \\
$z$-score & 0.722 & 0.468 & 0.213 & 0.331 & 0.766 \\
LOF & 0.723 & 0.556 & 0.168 & 0.302 & 0.785 \\
One-Class SVM & 0.732 & 0.712 & 0.077 & 0.277 & 0.771 \\
PCA-recon & 0.731 & 0.707 & 0.081 & 0.276 & 0.794 \\
HBOS & 0.603 & 0.306 & 0.457 & 0.239 & 0.674 \\
Isolation Forest & 0.470 & 0.201 & 0.641 & 0.174 & 0.664 \\
COPOD & 0.044 & 0.024 & 0.977 & $-$0.004 & 0.548 \\
ECOD & 0.027 & 0.022 & 0.986 & $-$0.029 & 0.493 \\
\bottomrule
\end{tabular}
\end{table}

\section{Drift and Retraining}
\label{sec:drift}

\paragraph{Does normal operation drift?} Streaming SWaT's genuine normal week
through a static Isolation Forest fitted on its first 14 hours, the share of
normal readings that raise an alarm, after a transient spike in the first day,
stays near the design rate of 1\% for about two days and then jumps, at roughly
hour 48, to $40$--$94\%$ for the rest of the week (mean $12\%$ over the first 24
streamed hours, $76\%$ over the last 24; Fig.~\ref{fig:drift}, left). The change
is a step, consistent with a change in operating regime, rather than a slow trend. The
Elliptic Envelope fitted on the same window does not show this: its false-alarm
rate over the normal week is $0.9\%$. Drift is therefore real, but whether it
matters depends on the detector. Isolation Forest operates on all 120 raw-scale
features, whereas the Elliptic Envelope operates on ten whitened principal
components; a plausible, untested explanation is that slow level changes in
individual sensors fall mostly outside those components.

\paragraph{Retraining budget.} A retrain \emph{attempt} is a batch on which the
policy fires; it becomes a \emph{refit} only if the selected five-hour window holds
more than 1{,}000 rows, otherwise the current model is kept. To compare policies at
equal cost, the periodic policy was run every $P\in\{1,2,3,4,6,12,24,48\}$ batches,
which traces how detection skill grows with the number of refits; the KS-triggered
policy is then read against that curve (Table~\ref{tab:drift}, Fig.~\ref{fig:drift}).

\paragraph{With labelled retraining data.} When retraining rows are selected with
ground-truth labels (as in the initial analysis), Isolation Forest's skill over the
attack period rises steadily with the number of refits: MCC $0.167$ static (73\% of
normal rows flagged), $0.276$ with 5 refits, $0.380$ with 40, $0.456$ with 80, $0.485$
with 120 and $0.537$ when refitting after every batch (239). The KS-triggered policy
reached $0.488$ with 74 refits (79 attempts), above the periodic curve at 80 refits
and level with periodic retraining at 120 refits, so triggering buys a modest saving in
the number of refits. For the Elliptic Envelope, no budget changes the outcome: MCC
$0.707$--$0.731$ at every periodic budget and $0.709$ for the KS policy, against $0.723$
for the static model.

\paragraph{Without labels, retraining can fail.} A deployed detector does not know
which recent rows are normal. When Isolation Forest retrains on the rows it did
\emph{not} flag, it enters a self-reinforcing failure: once it flags most rows, too few
remain to refit on, attempts stop turning into refits (193 attempts gave 52 refits
under the KS policy; 61 attempts gave 10 refits at $P{=}4$), and the detector ends up
flagging everything (MCC $0.000$--$0.033$ at every budget except $P{=}2$, $0.286$).
Retraining on all recent rows avoids that collapse but absorbs the attacks into the
model of normal behaviour: recall falls to $0.27$--$0.51$ and MCC to $0.074$--$0.303$
whenever refits are frequent. The Elliptic Envelope tolerates predicted-normal selection
at every budget (MCC $0.683$--$0.748$). With all-rows retraining it loses skill only when
refits are frequent (MCC $0.332$--$0.530$ with 40 or more refits, $0.683$--$0.717$ with
20 or fewer), because each frequent refit takes in a larger share of the ongoing attacks.
The benefit of adaptation reported in the initial analysis therefore depended on labels
that a deployment would not have.

\paragraph{Cost.} On the host, an Isolation Forest refit on five hours of data took
about $0.08$\,s and an Elliptic Envelope refit about $1.1$\,s. The KS check (120
features against a reference of up to 50{,}000 rows) cost $51$--$117$\,s of CPU over the
244 batches in the label-selected and all-rows runs ($0.2$--$0.5$\,s per batch) and
$483$\,s in the collapsing predicted-normal run. For Isolation Forest the checks alone
cost five times more than refitting after every batch ($96$\,s against $19$\,s, oracle
selection), and every-batch refitting gave the best detection, so a drift trigger does not
pay. A trigger can pay only when refits are expensive relative to the check, as for the
Elliptic Envelope, which in this stream did not need retraining at all.

\begin{table}[t]
\centering\small
\caption{Retraining budget on SWaT: MCC over the attack period (point-wise) for each
detector and retraining-data selection, at every periodic budget $P$ (refit attempt
every $P$ hourly batches; 244 batches in total) and for the KS-triggered policy. The
number in brackets is the number of successful refits. Static: Isolation Forest
$0.167$, Elliptic Envelope $0.723$.}
\label{tab:drift}
\begin{tabular}{llcccccccc}
\toprule
 & Selection & $P{=}1$ & $P{=}2$ & $P{=}3$ & $P{=}6$ & $P{=}12$ & $P{=}24$ & $P{=}48$ & KS \\
\midrule
IF & oracle & 0.537 (239) & 0.485 (120) & 0.456 (80) & 0.380 (40) & 0.325 (20) & 0.254 (10) & 0.276 (5) & 0.488 (74) \\
 & pred.-normal & 0.001 (52) & 0.286 (105) & 0.004 (25) & 0.000 (7) & 0.000 (4) & 0.033 (4) & 0.005 (1) & 0.025 (52) \\
 & all rows & 0.226 (244) & 0.147 (122) & 0.303 (81) & 0.124 (40) & 0.301 (20) & 0.225 (10) & 0.200 (5) & 0.188 (79) \\
EE & oracle & 0.711 (239) & 0.707 (120) & 0.717 (80) & 0.726 (40) & 0.719 (20) & 0.724 (10) & 0.731 (5) & 0.709 (74) \\
 & pred.-normal & 0.683 (238) & 0.729 (120) & 0.734 (80) & 0.711 (39) & 0.738 (20) & 0.748 (10) & 0.740 (5) & 0.725 (73) \\
 & all rows & 0.335 (244) & 0.405 (122) & 0.418 (81) & 0.423 (40) & 0.717 (20) & 0.701 (10) & 0.683 (5) & 0.332 (79) \\
\bottomrule
\end{tabular}
\end{table}

\begin{figure}[t]
\centering
\includegraphics[width=\textwidth]{drift_policies.png}
\caption{Left: false-alarm rate of a static Isolation Forest over SWaT's genuine normal
week (threshold set to flag 1\% of its training window). Middle and right: detection
skill (MCC) over the attack period against the number of successful refits; lines are
periodic retraining at different budgets, stars the KS-triggered policy, the dashed line
the static model. Colours give the retraining-data selection.}
\label{fig:drift}
\end{figure}

\paragraph{A note on the Elliptic Envelope's instability.} In this streaming
setting the static Elliptic Envelope, fitted on 14 hours, reaches MCC $0.723$,
yet in Table~\ref{tab:swat}, fitted on the full normal week with its threshold
set on the week's last 15\%, its label-free MCC is $0.008$. The two runs differ
only in the training window and threshold-calibration data. One consistent reading
is that drift within the normal week inflates the end-of-week scores used for
calibration, producing a threshold too high to fire. Either way, a detector
whose performance swings this much with the calibration window is fragile in
deployment, a property a single benchmark number hides.

\section{Resource Cost per Reading}
\label{sec:edge}

Table~\ref{tab:edge} and Fig.~\ref{fig:edge} report the cost of each detector
fitted on SWaT-shaped data (120 features; 50{,}000 normal rows) inside the
one-CPU, 1\,GB ARM64 container. Single-reading latency spans more than five orders of
magnitude. The $z$-score scores a reading in $0.002$\,ms; PCA reconstruction, HBOS,
Elliptic Envelope and One-Class SVM take $0.05$--$0.16$\,ms; LOF $0.89$\,ms;
Isolation Forest $2.3$\,ms, dominated by per-call overhead across its 100 trees,
since its batch throughput is about $850{,}000$ rows/s. ECOD and COPOD, as
implemented in PyOD, keep their training data and rescore against it on every
call: $698$ and $739$\,ms per reading and $235$\,MB models. Above a baseline of
about $444$\,MB (Python, libraries and data), the detectors themselves needed
$71$--$228$\,MB, except ECOD and COPOD ($516$ and $527$\,MB). The container's memory
use reached its 1\,GB limit only during the ECOD and COPOD runs: the kernel recorded
3{,}127 reclaim events at the limit and no out-of-memory kills, so the cap was
enforced and both completed. On SKAB-shaped data (40 features) the pattern is the
same at smaller scale (Isolation Forest $2.3$\,ms; ECOD and COPOD $118$--$125$\,ms,
39\,MB models, about $260$\,MB model memory). The LSTM autoencoder takes $0.78$\,ms
per 30-second window with 58k parameters, but needs a deep-learning runtime (peak
process memory $335$\,MB for inference alone). Recomputing the rolling features for
one new reading from a 30-row buffer with the pipeline's pandas code costs
$3.2$\,ms, more than any detector except ECOD and COPOD, so in a streaming
deployment incremental feature updates would matter more than the choice among the
fast detectors. These are lower bounds for a physical device. (An earlier
measurement in forked child processes over-counted memory, because a forked child's
resident set includes pages shared with its parent; it is not used here.)

\begin{table}[t]
\centering\small
\caption{Cost per detector in a Docker container capped at one CPU and 1\,GB RAM
(linux/arm64 on an Apple M4 host; single BLAS thread; one spawned process per
detector). SWaT-shaped input, 120 features, 50{,}000 training rows. Latency: one
reading per call, 1{,}000 calls after warm-up. Model memory: peak resident-memory
increase above the 444\,MB baseline during fitting and scoring.}
\label{tab:edge}
\begin{tabular}{lrrrrrrr}
\toprule
Detector & Fit (s) & Load (ms) & p50 (ms) & p95 (ms) & Batch rows/s & Size (KB) & Model mem.\ (MB) \\
\midrule
$z$-score & 0.006 & 0.12 & 0.002 & 0.002 & 7{,}188{,}644 & 2 & 107 \\
HBOS & 1.10 & 0.24 & 0.050 & 0.053 & 812{,}430 & 802 & 164 \\
PCA-recon & 0.019 & 0.18 & 0.058 & 0.063 & 1{,}242{,}616 & 52 & 71 \\
Elliptic Envelope & 4.22 & 0.31 & 0.151 & 0.161 & 5{,}866{,}393 & 503 & 89 \\
One-Class SVM & 3.10 & 2.78 & 0.158 & 0.172 & 10{,}049 & 1{,}942 & 228 \\
LOF & 2.55 & 1.88 & 0.892 & 0.934 & 7{,}745 & 22{,}188 & 128 \\
Isolation Forest & 0.072 & 7.46 & 2.255 & 2.353 & 853{,}367 & 1{,}149 & 85 \\
ECOD & 0.83 & 196.5 & 698.2 & 956.7 & 8{,}960 & 235{,}157 & 516 \\
COPOD & 1.11 & 222.6 & 739.3 & 833.9 & 11{,}203 & 235{,}157 & 527 \\
\bottomrule
\end{tabular}
\end{table}

\begin{figure}[t]
\centering
\includegraphics[width=0.75\textwidth]{edge_cost.png}
\caption{Single-reading latency against serialised model size (both log scale) in the
capped container, for SKAB-shaped (40 features) and SWaT-shaped (120 features) input.}
\label{fig:edge}
\end{figure}

\section{Discussion}
\label{sec:discussion}

Table~\ref{tab:synth} condenses what each protocol choice did to the measured outcome
and what it implies for someone selecting or deploying a detector.

\begin{table}[t]
\centering\small
\caption{Summary: how each evaluation choice changed the measured outcome in this study,
and the practical implication.}
\label{tab:synth}
\begin{tabular}{p{0.17\textwidth}p{0.43\textwidth}p{0.32\textwidth}}
\toprule
Choice & Evidence in this study & Implication \\
\midrule
Sensor set & Adding one drifting SWaT analyser (AIT201): $z$-score false alarms $1.4\%\to42\%$. Removing three drifting WADI analysers: false alarms $56\%\to1\%$ for six detectors & Check sensor ranges against training data; prefer detectors that bound each sensor's influence when analysers can drift \\
Training/test split & SWaT re-upload layout vs standard split: AUROC $0.84$--$0.93$ vs $0.83$--$0.86$; HBOS first vs last. SKAB pooled vs per-run: chance-level vs run AUROC up to $0.79$, ranking inverted & State the time span and source of training and test data; use a benchmark's reference protocol \\
Threshold rule & Shared contamination vs per-model calibration: $z$-score MCC $0.18$ vs $0.84$. Label-free vs label-calibrated on SWaT: $z$-score MCC $0.22$ vs $0.82$, One-Class SVM $0.18$ vs $0.82$ & Report a label-free threshold as the deployable result; treat label-calibrated results as an optimistic reference \\
Retraining data & Oracle labels: Isolation Forest MCC $0.17\to0.54$. Predicted-normal: collapse to flag-everything. All rows: attacks absorbed, MCC $\le0.30$ & Do not claim adaptation gains from label-selected retraining; vet retraining data or use slowly changing detectors \\
Retraining trigger & KS trigger: $0.488$ with 74 refits vs periodic $0.456$ (80), $0.485$ (120); KS checks $96$\,s vs every-batch refits $19$\,s & Compare triggers against a budget curve, including the cost of the check \\
Cost measurement & Batch-averaged vs single-reading latency for Isolation Forest: $0.0019$ vs $2.3$\,ms; feature update $3.2$\,ms & Report single-reading latency and memory above a stated baseline \\
\bottomrule
\end{tabular}
\end{table}

\paragraph{Ranking quality is not the bottleneck; thresholds and sensors are.}
On SWaT, ten detectors that span statistical, distance, density, tree, copula and
recurrent families are statistically indistinguishable in AUROC, and on WADI they become
so once three drifting analysers are removed. This agrees with large benchmark studies
that find no consistent winner and strong simple baselines
\citep{audibert2022,sarfraz2024}. What separates detectors in use is whether a threshold
set on normal data still works later. On both water plants the decisive factor was the
chemical analysers, whose readings left their training range between the normal and the
attack recordings: detectors that combine all features through distances or
reconstructions inherited that shift and raised alarms on roughly half or more of normal
readings, while detectors that bound each feature's influence (COPOD, ECOD, HBOS,
Isolation Forest) did not. On HAI, where no comparable shift occurred, thresholds
transferred and the distance-based detectors led. Which detector is best is therefore a
property of the plant, its sensors and the protocol together. For a deployment this
suggests two practical checks before choosing a detector: compare recent sensor ranges
with the training ranges, and evaluate thresholds on normal data recorded after the
training period.

\paragraph{Why the initial conclusions changed.} Three protocol choices in the
initial analysis each favoured a particular answer. Training on the attack period's
own normal rows removed the normal-week to attack-week shift from the evaluation.
A contamination value tuned on Isolation Forest and shared by all detectors made
detectors with differently scaled scores look poor. Retraining on
ground-truth-normal rows made adaptation look free of risk. None of these
choices is unusual, and each was documented in the code. Their effect became
visible only when the same detectors were rerun under alternative protocols.
Studies that compare detectors should therefore report at least one alternative
threshold rule and state where their training and test data come from in time.

\paragraph{What adaptation needs.} Normal operation does drift, and some
detectors are much more sensitive to it than others. Where retraining helps, it
helps only if the retraining data are clean; with label-free selection Isolation
Forest either collapsed into flagging everything or learned the attacks as
normal. In practice this argues for one of three things: retraining data vetted
by an operator; a slow-changing detector (here, the Elliptic Envelope on whitened
components), whose static form was as good as any adaptive variant; or adaptation
of the threshold rather than the model. For cheap models such as Isolation Forest, a
drift test can cost more than the retrain it is meant to save.

\paragraph{Cost.} Single-reading latencies span almost six orders of magnitude, from
microseconds ($z$-score) to about 700\,ms (ECOD and COPOD in PyOD on SWaT, which store and
rescore against their training data), and model sizes from 1\,KB to 235\,MB. The
batch-amortised figures used in the initial analysis (around $0.002$\,ms per row for
Isolation Forest) understate the real single-reading cost (about $2.3$--$2.8$\,ms)
by three orders of magnitude, because per-call overhead dominates. For a monitor that
scores one reading per second, all detectors except possibly ECOD and COPOD fit
comfortably; the deciding constraints are memory and threshold stability, not
speed.

\section{Threats to Validity and Limitations}
\label{sec:limits}

\begin{itemize}
\item \textbf{Dataset versions.} SWaT and WADI results use the official iTrust releases
(SWaT A1 \& A2, December 2015; WADI A2, 2019). Later SWaT and WADI campaigns, with
different attacks and plant configurations, were not evaluated. The drift study and the
cost benchmark used SWaT; the cost benchmark used the re-upload's 24 sensors, which does
not affect timing materially.
\item \textbf{Process coverage.} SKAB and SWaT are water processes; HAI adds
turbine and boiler processes but through a hardware-in-the-loop simulator. Rotating
machinery under real wear, electrical faults and naturally occurring (rather than
induced or scripted) anomalies are not covered; WADI \citep{ahmed2017} and run-to-failure
machinery data are natural next steps. The drift and cost studies were run on SWaT
only.
\item \textbf{Detector configurations.} Hyperparameters other than thresholds were
fixed at common defaults; One-Class SVM and LOF were trained on a 20{,}000-row
subsample on SWaT and HAI, and ECOD and COPOD on a 100{,}000-row subsample on HAI; the LSTM autoencoder is small and was run with one seed on
SWaT. Larger or tuned deep models were not tested.
\item \textbf{Metrics.} Event recall and median delay are reported, but
range-aware metrics such as VUS-PR \citep{paparrizos2022} were not computed.
Bootstrap blocks (files, hours) reduce but do not remove dependence between blocks.
\item \textbf{Drift design.} One stream, one window length (five hours), one KS
threshold rule and two detectors were studied; ADWIN-style detectors and threshold-only
adaptation were not compared. Policies are compared without confidence intervals
because the stream is a single realisation.
\item \textbf{Resource measurements.} A Docker container capped at one CPU and
1\,GB on an Apple M4 host is not a physical edge device: slower cores, thermal
limits and real I/O are absent, so latencies are lower bounds. Model memory is the peak
increase above a measured baseline and includes transient scoring buffers. Energy was not measured.
\item \textbf{SKAB protocols.} The SKAB reference protocol trains on rows of which
2.1\% are labelled anomalous, and its label-free threshold is set from each
run's own training scores; results are not directly comparable with the SKAB
leaderboard, whose features, thresholds and post-processing differ.
\item \textbf{Search scope.} The literature search used a general web search and
publisher pages rather than database-level systematic queries; statements that a
comparison was not found are limited to that search.
\end{itemize}

\section{Conclusion}
\label{sec:conclusion}

Evaluating ten lightweight and deep detectors on SKAB, SWaT, WADI and HAI under
documented, standard protocols answered the four research questions as follows.
(RQ1) On SWaT the detectors rank attacks equally well (AUROC $0.81$--$0.86$, overlapping
intervals), and on WADI they do so once three drifting analysers are removed; on HAI
they separate, with PCA reconstruction and LOF leading ($0.93$, $0.92$) and Elliptic
Envelope trailing ($0.66$); on SKAB with pooled training, none is reliably better than
chance. (RQ2) Protocol choices change the conclusions more than detector choice does:
the SWaT split definition reorders the detectors, a shared contamination value had made
the $z$-score look useless (MCC $0.18$ against $0.84$ with its own threshold), SKAB's
per-run protocol inverts the pooled ranking, and a single analyser decides whether
label-free thresholds survive (the $z$-score's false alarms $1.4\%$ without AIT201,
$42\%$ with it). With label-free thresholds on the official SWaT data, COPOD, Isolation
Forest and HBOS were the most dependable (MCC $0.46$--$0.63$). (RQ3) Normal operation drifts enough to break a static Isolation
Forest; retraining restores it only when the retraining data are known to be
normal, drift-triggered retraining saved some refits but
cost more in drift checks than it saved, and label-free retraining failed. (RQ4) Single-reading costs range from
microseconds to about 700\,ms and model sizes from kilobytes to hundreds of megabytes
under a one-CPU, 1\,GB cap; the per-reading feature update (3.2\,ms) costs more than
most detectors.

The initial analysis reported dataset-specific winners, a large
benefit from drift adaptation, and sub-millisecond inference; none of these
survived re-examination. What did survive is a clearer and more useful message for
practitioners: choose the detector for its threshold stability and retraining
needs rather than its benchmark rank, and report the protocol in enough detail
that a reader can tell which of these the numbers measure. Next steps are later SWaT
and WADI campaigns, data from rotating machinery under real wear, drift-aware
handling of analyser channels, threshold-only adaptation, and measurement on physical
edge hardware.

\section*{Data and Code Availability}
Code, configuration and all result files are available in the project repository
(\texttt{experiments/revision/}, \texttt{results/revised/}). SKAB and HAI are public
\citep{skab,shin2021}. SWaT and WADI must be requested from iTrust, SUTD, and are not
redistributed; the comparison with the public SWaT re-upload is documented in the code. Software used:
scikit-learn \citep{pedregosa2011}, PyOD \citep{zhao2019} and PyTorch.

\section*{Disclosure of Interests}
The authors have no competing interests to declare.

\begin{thebibliography}{99}
\small

\bibitem[Ahmad et~al.(2017)]{ahmad2017} Ahmad, S., Lavin, A., Purdy, S., \& Agha, Z. (2017). Unsupervised real-time anomaly detection for streaming data. \emph{Neurocomputing}, 262, 134--147. doi:10.1016/j.neucom.2017.04.070

\bibitem[Ahmed et~al.(2017)]{ahmed2017} Ahmed, C.~M., Palleti, V.~R., \& Mathur, A.~P. (2017). WADI: A water distribution testbed for research in the design of secure cyber physical systems. In \emph{Proc.\ 3rd Int.\ Workshop on Cyber-Physical Systems for Smart Water Networks (CySWater)} (pp.\ 25--28). doi:10.1145/3055366.3055375

\bibitem[Audibert et~al.(2020)]{audibert2020} Audibert, J., Michiardi, P., Guyard, F., Marti, S., \& Zuluaga, M.~A. (2020). USAD: Unsupervised anomaly detection on multivariate time series. In \emph{Proc.\ 26th ACM SIGKDD Int.\ Conf.\ on Knowledge Discovery \& Data Mining}. doi:10.1145/3394486.3403392

\bibitem[Audibert et~al.(2022)]{audibert2022} Audibert, J., Michiardi, P., Guyard, F., Marti, S., \& Zuluaga, M.~A. (2022). Do deep neural networks contribute to multivariate time series anomaly detection? \emph{Pattern Recognition}, 132, 108945. doi:10.1016/j.patcog.2022.108945

\bibitem[Banbury et~al.(2021)]{banbury2021} Banbury, C., et~al. (2021). MLPerf Tiny benchmark. In \emph{Proc.\ NeurIPS Datasets and Benchmarks Track}. arXiv:2106.07597

\bibitem[Bifet \& Gavald\`a(2007)]{bifet2007} Bifet, A., \& Gavald\`a, R. (2007). Learning from time-changing data with adaptive windowing. In \emph{Proc.\ SIAM Int.\ Conf.\ on Data Mining} (pp.\ 443--448). doi:10.1137/1.9781611972771.42

\bibitem[Breunig et~al.(2000)]{breunig2000} Breunig, M.~M., Kriegel, H.-P., Ng, R.~T., \& Sander, J. (2000). LOF: Identifying density-based local outliers. In \emph{Proc.\ ACM SIGMOD Int.\ Conf.\ on Management of Data} (pp.\ 93--104). doi:10.1145/342009.335388

\bibitem[Chandola et~al.(2009)]{chandola2009} Chandola, V., Banerjee, A., \& Kumar, V. (2009). Anomaly detection: A survey. \emph{ACM Computing Surveys}, 41(3), 1--58. doi:10.1145/1541880.1541882

\bibitem[Chicco \& Jurman(2020)]{chicco2020} Chicco, D., \& Jurman, G. (2020). The advantages of the Matthews correlation coefficient (MCC) over F1 score and accuracy in binary classification evaluation. \emph{BMC Genomics}, 21, 6. doi:10.1186/s12864-019-6413-7

\bibitem[Dasari(2026)]{dasari2026} Dasari, S. (2026). When to retrain: An empirical study of retraining policies for streaming ML under concept drift, budget, and latency constraints. arXiv preprint arXiv:2608.19488.

\bibitem[Deng \& Hooi(2021)]{deng2021} Deng, A., \& Hooi, B. (2021). Graph neural network-based anomaly detection in multivariate time series. In \emph{Proc.\ AAAI Conf.\ on Artificial Intelligence}, 35(5), 4027--4035. doi:10.1609/aaai.v35i5.16523

\bibitem[Gama et~al.(2014)]{gama2014} Gama, J., \v{Z}liobait\.{e}, I., Bifet, A., Pechenizkiy, M., \& Bouchachia, A. (2014). A survey on concept drift adaptation. \emph{ACM Computing Surveys}, 46(4), 1--37. doi:10.1145/2523813

\bibitem[Garg et~al.(2021)]{garg2021} Garg, A., Zhang, W., Samaran, J., Savitha, R., \& Foo, C.-S. (2021). An evaluation of anomaly detection and diagnosis in multivariate time series. \emph{IEEE Transactions on Neural Networks and Learning Systems}. doi:10.1109/TNNLS.2021.3105827

\bibitem[Goh et~al.(2016)]{goh2016} Goh, J., Adepu, S., Junejo, K.~N., \& Mathur, A. (2016). A dataset to support research in the design of secure water treatment systems. In \emph{Proc.\ 11th Int.\ Conf.\ on Critical Information Infrastructures Security (CRITIS)} (pp.\ 88--99). doi:10.1007/978-3-319-71368-7\_8

\bibitem[Goldstein \& Uchida(2016)]{goldstein2016} Goldstein, M., \& Uchida, S. (2016). A comparative evaluation of unsupervised anomaly detection algorithms for multivariate data. \emph{PLoS ONE}, 11(4), e0152173. doi:10.1371/journal.pone.0152173

\bibitem[Han et~al.(2022)]{han2022} Han, S., Hu, X., Huang, H., Jiang, M., \& Zhao, Y. (2022). ADBench: Anomaly detection benchmark. In \emph{Advances in Neural Information Processing Systems 35 (Datasets and Benchmarks Track)}. arXiv:2206.09426

\bibitem[Hundman et~al.(2018)]{hundman2018} Hundman, K., Constantinou, V., Laporte, C., Colwell, I., \& Soderstrom, T. (2018). Detecting spacecraft anomalies using LSTMs and nonparametric dynamic thresholding. In \emph{Proc.\ 24th ACM SIGKDD Int.\ Conf.\ on Knowledge Discovery \& Data Mining}. doi:10.1145/3219819.3219845

\bibitem[Katser \& Kozitsin(2020)]{skab} Katser, I.~D., \& Kozitsin, V.~O. (2020). \emph{Skoltech Anomaly Benchmark (SKAB)} [Data set]. Kaggle. doi:10.34740/KAGGLE/DSV/1693952

\bibitem[Kim et~al.(2022)]{kim2022} Kim, S., Choi, K., Choi, H.-S., Lee, B., \& Yoon, S. (2022). Towards a rigorous evaluation of time-series anomaly detection. In \emph{Proc.\ AAAI Conf.\ on Artificial Intelligence}, 36(7), 7194--7201. doi:10.1609/aaai.v36i7.20680

\bibitem[Lee(2017)]{lee2017} Lee, W.-J. (2017). Anomaly detection and severity prediction of air leakage in train braking pipes. \emph{International Journal of Prognostics and Health Management}, 8(3). doi:10.36001/ijphm.2017.v8i3.2662

\bibitem[Li et~al.(2019)]{li2019madgan} Li, D., Chen, D., Jin, B., Shi, L., Goh, J., \& Ng, S.-K. (2019). MAD-GAN: Multivariate anomaly detection for time series data with generative adversarial networks. In \emph{Proc.\ Int.\ Conf.\ on Artificial Neural Networks (ICANN)}. arXiv:1809.04758

\bibitem[Li et~al.(2023)]{li2022ecod} Li, Z., Zhao, Y., Hu, X., Botta, N., Ionescu, C., \& Chen, G.~H. (2023). ECOD: Unsupervised outlier detection using empirical cumulative distribution functions. \emph{IEEE Transactions on Knowledge and Data Engineering}, 35(12), 12181--12193. doi:10.1109/TKDE.2022.3159580

\bibitem[Liu \& Paparrizos(2024)]{liu2024} Liu, Q., \& Paparrizos, J. (2024). The elephant in the room: Towards a reliable time-series anomaly detection benchmark. In \emph{Advances in Neural Information Processing Systems 37 (Datasets and Benchmarks Track)} (pp.\ 108231--108261).

\bibitem[Liu et~al.(2008)]{liu2008} Liu, F.~T., Ting, K.~M., \& Zhou, Z.-H. (2008). Isolation forest. In \emph{Proc.\ 8th IEEE Int.\ Conf.\ on Data Mining} (pp.\ 413--422). doi:10.1109/ICDM.2008.17

\bibitem[Lu et~al.(2019)]{lu2018} Lu, J., Liu, A., Dong, F., Gu, F., Gama, J., \& Zhang, G. (2019). Learning under concept drift: A review. \emph{IEEE Transactions on Knowledge and Data Engineering}, 31(12), 2346--2363. doi:10.1109/TKDE.2018.2876857

\bibitem[Malhotra et~al.(2016)]{malhotra2016} Malhotra, P., Ramakrishnan, A., Anand, G., Vig, L., Agarwal, P., \& Shroff, G. (2016). LSTM-based encoder-decoder for multi-sensor anomaly detection. \emph{ICML Anomaly Detection Workshop}. arXiv:1607.00148

\bibitem[Massey(1951)]{massey1951} Massey, F.~J. (1951). The Kolmogorov-Smirnov test for goodness of fit. \emph{Journal of the American Statistical Association}, 46(253), 68--78. doi:10.1080/01621459.1951.10500769

\bibitem[Mathur \& Tippenhauer(2016)]{mathur2016} Mathur, A.~P., \& Tippenhauer, N.~O. (2016). SWaT: A water treatment testbed for research and training on ICS security. In \emph{Proc.\ Int.\ Workshop on Cyber-physical Systems for Smart Water Networks (CySWater)} (pp.\ 31--36). doi:10.1109/CySWater.2016.7469060

\bibitem[Mukherjee(2021)]{mukherjee2021} Mukherjee, R. (2021). A modified energy statistic for unsupervised anomaly detection. \emph{International Journal of Prognostics and Health Management}, 12(1). doi:10.36001/ijphm.2021.v12i1.1323

\bibitem[Paparrizos et~al.(2022)]{paparrizos2022} Paparrizos, J., Boniol, P., Palpanas, T., Tsay, R.~S., Elmore, A., \& Franklin, M.~J. (2022). Volume under the surface: A new accuracy evaluation measure for time-series anomaly detection. \emph{Proceedings of the VLDB Endowment}, 15(11), 2774--2787. doi:10.14778/3551793.3551830

\bibitem[Pedregosa et~al.(2011)]{pedregosa2011} Pedregosa, F., et~al. (2011). Scikit-learn: Machine learning in Python. \emph{Journal of Machine Learning Research}, 12, 2825--2830.

\bibitem[Poenaru-Olaru et~al.(2024)]{poenaru2024} Poenaru-Olaru, L., Karpova, N., Cruz, L., Rellermeyer, J.~S., \& van Deursen, A. (2024). Is your anomaly detector ready for change? Adapting AIOps solutions to the real world. In \emph{Proc.\ IEEE/ACM 3rd Int.\ Conf.\ on AI Engineering -- Software Engineering for AI (CAIN)} (pp.\ 222--233). doi:10.1145/3644815.3644961

\bibitem[Rabanser et~al.(2019)]{rabanser2019} Rabanser, S., G\"unnemann, S., \& Lipton, Z. (2019). Failing loudly: An empirical study of methods for detecting dataset shift. In \emph{Advances in Neural Information Processing Systems 32}. arXiv:1810.11953

\bibitem[Rousseeuw \& Van~Driessen(1999)]{rousseeuw1999} Rousseeuw, P.~J., \& Van~Driessen, K. (1999). A fast algorithm for the minimum covariance determinant estimator. \emph{Technometrics}, 41(3), 212--223. doi:10.1080/00401706.1999.10485670

\bibitem[Sarfraz et~al.(2024)]{sarfraz2024} Sarfraz, M.~S., Chen, M.-Y., Layer, L., Peng, K., \& Koulakis, M. (2024). Position: Quo vadis, unsupervised time series anomaly detection? In \emph{Proc.\ 41st Int.\ Conf.\ on Machine Learning}, PMLR 235, 43461--43476.

\bibitem[Schmidl et~al.(2022)]{schmidl2022} Schmidl, S., Wenig, P., \& Papenbrock, T. (2022). Anomaly detection in time series: A comprehensive evaluation. \emph{Proceedings of the VLDB Endowment}, 15(9), 1779--1797. doi:10.14778/3538598.3538602

\bibitem[Sch\"olkopf et~al.(2001)]{scholkopf2001} Sch\"olkopf, B., Platt, J.~C., Shawe-Taylor, J., Smola, A.~J., \& Williamson, R.~C. (2001). Estimating the support of a high-dimensional distribution. \emph{Neural Computation}, 13(7), 1443--1471. doi:10.1162/089976601750264965

\bibitem[Shin et~al.(2020)]{shin2020} Shin, H.-K., Lee, W., Yun, J.-H., \& Kim, H. (2020). HAI 1.0: HIL-based augmented ICS security dataset. In \emph{Proc.\ 13th USENIX Workshop on Cyber Security Experimentation and Test (CSET)}.

\bibitem[Shin et~al.(2021)]{shin2021} Shin, H.-K., Lee, W., Yun, J.-H., \& Min, B.-G. (2021). Two ICS security datasets and anomaly detection contest on the HIL-based augmented ICS testbed. In \emph{Proc.\ Cyber Security Experimentation and Test Workshop (CSET '21)}.

\bibitem[Shyu et~al.(2003)]{shyu2003} Shyu, M.-L., Chen, S.-C., Sarinnapakorn, K., \& Chang, L. (2003). A novel anomaly detection scheme based on principal component classifier. In \emph{Proc.\ IEEE ICDM Workshop on Foundations and New Directions of Data Mining} (pp.\ 172--179).

\bibitem[Tatbul et~al.(2018)]{tatbul2018} Tatbul, N., Lee, T.~J., Zdonik, S., Alam, M., \& Gottschlich, J. (2018). Precision and recall for time series. In \emph{Advances in Neural Information Processing Systems 31}. arXiv:1803.03639

\bibitem[Tuli et~al.(2022)]{tuli2022} Tuli, S., Casale, G., \& Jennings, N.~R. (2022). TranAD: Deep transformer networks for anomaly detection in multivariate time series data. \emph{Proceedings of the VLDB Endowment}, 15(6), 1201--1214. doi:10.14778/3514061.3514067

\bibitem[Weiss \& Brundage(2021)]{weiss2021} Weiss, B., \& Brundage, M. (2021). Measurement and evaluation for prognostics and health management (PHM) for manufacturing operations: Summary of an interactive workshop highlighting PHM trends. \emph{International Journal of Prognostics and Health Management}, 12(1). doi:10.36001/ijphm.2021.v12i1.2653

\bibitem[Wu \& Keogh(2023)]{wu2023} Wu, R., \& Keogh, E. (2023). Current time series anomaly detection benchmarks are flawed and are creating the illusion of progress. \emph{IEEE Transactions on Knowledge and Data Engineering}, 35(3), 2421--2429. doi:10.1109/TKDE.2021.3112126

\bibitem[Zhao et~al.(2019)]{zhao2019} Zhao, Y., Nasrullah, Z., \& Li, Z. (2019). PyOD: A Python toolbox for scalable outlier detection. \emph{Journal of Machine Learning Research}, 20(96), 1--7.

\end{thebibliography}

\end{document}
